\documentclass[12pt]{article}
\usepackage[utf8]{inputenc}
\usepackage[T1]{fontenc}
\usepackage{hyperref}
\title{Kanashi Language Resources for K-12 Curriculum Development}
\author{Compiled by Claude Sonnet 5.5 (claude-sonnet-5-5) for Andr\'as Kornai}
\date{}
\begin{document}
\maketitle

\section{Introduction}
Kanashi (ISO 639-3 xns; Glottocode kana1283), also called Malani or Kanasi, is a Sino-Tibetan language of the West Himalayish branch. It is spoken in one village, Malana in Kullu district, Himachal Pradesh, India. Speaker estimates vary: Wikipedia gives 1,700 (2015), and Saxena and Borin (2022) say about 2,000. The language is partly mutually intelligible with Kinnauri. It is surrounded by Indo-Aryan languages. It is not known to be taught in schools.

\section{Primary Educational Resources}
The search found no primers, textbooks or school programmes for Kanashi.

\section{Linguistic and Reference Materials}
\begin{itemize}
\item \emph{Synchronic and Diachronic Aspects of Kanashi}, ed. Anju Saxena and Lars Borin (De Gruyter Mouton, 2022). It is an open-access volume covering grammar, sound system, numerals, vocabulary and genealogy. Open-access PDF: \url{https://library.oapen.org/bitstream/20.500.12657/53703/1/9783110703245.pdf}
\item Publisher page: \url{https://www.degruyterbrill.com/isbn/9783110703245}
\item Kanashi language, Wikipedia overview: \url{https://en.wikipedia.org/wiki/Kanashi_language}
\item Malana village, Wikipedia: \url{https://en.wikipedia.org/wiki/Malana,_Himachal_Pradesh}
\end{itemize}

\section{Religious and Community Resources}
The search found no religious or community-organisation resources specific to Kanashi.

\section{Digital and Audio Resources}
No audio, video or software resources were found in these searches.

\section{Curriculum Development Context}
The language is confined to one village and has very little published documentation beyond the 2022 volume. Any curriculum would have to be developed with Malana community members, using Saxena and Borin as the linguistic base. No orthography or teaching materials were found, and no literacy tradition in Kanashi was identified.

\end{document}